\section{PCIe Traffic and Speculative Decoding}\label{sec:bytes}

The samplers record each GPU's PCIe receive rate every half second (Table~\ref{tab:pcie-per-token}). Summed over both
GPUs it peaks at \dv{ctl.pcie.E8.B.N32.gbs}\,GB/s in E8 and \dv{ctl.pcie.E13.A.N12.gbs}\,GB/s in the coding-agent
experiment, against a line rate of \dv{ctl.pcie.calib.link_gbs}\,GB/s per GPU and direction for a PCIe Gen5 x16 link.
In E8 the bytes received per output token fall as users are added, from \dv{ctl.pcie.E8.A.N1.mb_per_tok} to
\dv{ctl.pcie.E8.A.N64.mb_per_tok}\,MB with unique prompts and from \dv{ctl.pcie.E8.B.N1.mb_per_tok} to
\dv{ctl.pcie.E8.B.N64.mb_per_tok}\,MB with the cached prefix, which is what one expects if one read of a host expert
serves several requests of a batch. In the coding-agent experiment they do not fall
(\dv{ctl.pcie.E13.A.N1.mb_per_tok}\,MB at one stream, \dv{ctl.pcie.E13.A.N12.mb_per_tok}\,MB at twelve in its short-session
regime), and we cannot explain the difference with these data. The two experiments differ in reply length (400 tokens
against 256), host (pod~3 against pod~1) and how far past saturation their sweeps went, and the counter mixes prompt
reading, drafter reads and expert reads, which we did not attribute. The counter itself is only partly validated. On
pod~1 a calibration script found NVML reporting \dv{ctl.pcie.calib.gpu0.copy.ratio} and
\dv{ctl.pcie.calib.gpu1.copy.ratio} times the true bytes for copies on GPU~0 and~1, and
\dv{ctl.pcie.calib.gpu0.zerocopy.ratio} times for zero-copy reads on GPU~0. The zero-copy test on GPU~1 is invalid, and
pods~2 and~3 have no calibration, so these values are raw and descriptive.

DSpark~\citep{cheng2026dspark,leviathan2023speculative}, the model's drafter, proposes \dv{ctl.spec.k} draft tokens per
step, and on code prompts a request gains on average \dv{ctl.spec.E8.A.accept_len} tokens per verification step with
unique prompts and \dv{ctl.spec.E8.B.accept_len} with the cached prefix (\dv{ctl.spec.E8.A.accept_rate_pct}\% and
\dv{ctl.spec.E8.B.accept_rate_pct}\% of proposed tokens accepted), almost independent of the number of users (in A
\dv{ctl.spec.E8.A.N1.accept_len} with one user and \dv{ctl.spec.E8.A.N64.accept_len} with 64,
Table~\ref{tab:spec-acceptance}). The coding-agent regimes, pooled over their levels, range from
\dv{ctl.spec.E13.A.accept_len} for short independent prompts to \dv{ctl.spec.E13.C256k.accept_len} for independent
256K-token contexts. All prompts were code or coding-agent turns with thinking off, so these values say nothing about
prose, and with speculation always on, its share of the speeds reported in the paper is not separated from the
offload's (E7, which would switch it off, did not run).

% Generated by hostoffload.analysis.pipeline. Do not edit.
\begin{table}[t]
\centering
\small
\begin{fittable}{tab:pcie-per-token}
\begin{tabular}{lrrr}
\toprule
regime & $N$ & PCIe RX GB/s & MB per output token \\
\midrule
E8 A: 4K independent & 1 & 25.2 & 473 \\
 & 8 & 72.4 & 395 \\
 & 64 & 66.3 & 231 \\
E8 B: 16K cached prefix & 1 & 36.5 & 326 \\
 & 8 & 76.2 & 356 \\
 & 64 & 68.2 & 150 \\
E13 short: 4K independent & 1 & 15.1 & 365 \\
 & 8 & 75.3 & 373 \\
 & 12 & 91.3 & 599 \\
E13 shared 128K prefix & 1 & 26.5 & 290 \\
 & 8 & 78.6 & 324 \\
 & 12 & 90.7 & 515 \\
E13 shared 256K prefix & 1 & 27.9 & 312 \\
 & 8 & 76.4 & 330 \\
 & 12 & 88.7 & 497 \\
E13 own 128K context & 1 & 50.2 & 455 \\
 & 5 & 84.2 & 570 \\
E13 own 256K context & 1 & 49.1 & 464 \\
 & 3 & 59.6 & 325 \\
E13 agent turns & 1 & 49.0 & 412 \\
 & 8 & 81.6 & 386 \\
\bottomrule
\end{tabular}
\end{fittable}
\par\smallskip\parbox{0.95\linewidth}{\footnotesize\raggedright Calibration on pod 1 (\texttt{data/controlled/E8-main/pcie-calibration.json}): NVML over true bytes 1.11 for DMA copies and 1.22 for zero-copy kernel reads on GPU 0. The GPU 1 zero-copy test is invalid: its nominal 294 GB/s exceeds the 63 GB/s PCIe 5.0 x16 line rate and NVML saw almost no traffic, so the kernel did not read over PCIe. Values are not corrected.}
\caption{PCIe receive traffic summed over both GPUs (NVML, raw counter) per level, and megabytes received per output token completed in the level window. E8 ran on pod 1 (two repetitions for B and for A at $N=1$, pooled), E13 on pod 3 (one repetition). Rows: $N=1$, $N=8$ and the largest tested $N$; every level is in \texttt{values.tex} (\texttt{ctl.pcie.*}). Descriptive: the bytes include prefill and drafter traffic.}
\label{tab:pcie-per-token}
\end{table}

% Generated by hostoffload.analysis.pipeline. Do not edit.
\begin{table}[t]
\centering
\small
\begin{fittable}{tab:spec-acceptance}
\begin{tabular}{lrrrrr}
\toprule
regime & $N$ & drafts & accepted/draft & accept. length & accepted \% \\
\midrule
E8 A: 4K independent & 1 & 1,139 & 2.39 & 3.39 & 48 \\
 & 8 & 1,192 & 2.44 & 3.44 & 49 \\
 & 64 & 9,744 & 2.37 & 3.37 & 47 \\
 & all & 20,627 & 2.37 & 3.37 & 47 \\
E8 B: 16K cached prefix & 1 & 1,134 & 2.63 & 3.63 & 53 \\
 & 8 & 2,318 & 2.54 & 3.54 & 51 \\
 & 64 & 18,279 & 2.58 & 3.58 & 52 \\
 & all & 37,575 & 2.60 & 3.60 & 52 \\
E13 short: 4K independent & 1 & 450 & 2.57 & 3.57 & 51 \\
 & 8 & 1,903 & 2.37 & 3.37 & 47 \\
 & 12 & 2,852 & 2.37 & 3.37 & 47 \\
 & all & 9,991 & 2.37 & 3.37 & 47 \\
E13 shared 128K prefix & 1 & 472 & 2.39 & 3.39 & 48 \\
 & 8 & 1,677 & 2.82 & 3.82 & 56 \\
 & 12 & 2,760 & 2.48 & 3.48 & 50 \\
 & all & 9,336 & 2.60 & 3.60 & 52 \\
E13 shared 256K prefix & 1 & 417 & 2.85 & 3.85 & 57 \\
 & 8 & 1,763 & 2.64 & 3.64 & 53 \\
 & 12 & 2,718 & 2.54 & 3.54 & 51 \\
 & all & 9,315 & 2.62 & 3.62 & 52 \\
E13 own 128K context & 1 & 223 & 2.60 & 3.60 & 52 \\
 & 5 & 1,063 & 2.77 & 3.77 & 55 \\
 & all & 3,272 & 2.67 & 3.67 & 53 \\
E13 own 256K context & 1 & 194 & 3.15 & 4.15 & 63 \\
 & 3 & 653 & 2.67 & 3.67 & 53 \\
 & all & 1,275 & 2.77 & 3.77 & 55 \\
E13 agent turns & 1 & 302 & 2.98 & 3.98 & 60 \\
 & 8 & 2,594 & 2.71 & 3.71 & 54 \\
 & all & 7,726 & 2.74 & 3.74 & 55 \\
\bottomrule
\end{tabular}
\end{fittable}
\caption{DSpark acceptance per level from vLLM's counters: drafts (\texttt{vllm:spec\_decode\_num\_drafts}), accepted draft tokens per draft, acceptance length $1 + \text{accepted}/\text{drafts}$ (tokens a request gains per verification step) and the share of proposed draft tokens accepted. Rows: $N=1$, $N=8$, the largest tested $N$ and all levels pooled; E8 on pod 1 (repetitions pooled), E13 on pod 3.}
\label{tab:spec-acceptance}
\end{table}

